% Generated draft: requires mathematical and source review.
\documentclass[11pt]{article}
\usepackage[margin=25mm]{geometry}
\usepackage{fontspec}
\setmainfont{Latin Modern Roman}
\usepackage{amsmath,amssymb,mathtools,unicode-math}
\setmathfont{Latin Modern Math}
\usepackage{enumitem,needspace,xurl,hyperref}
\setlength{\parindent}{0pt}
\setlength{\parskip}{5pt}
\setlength{\emergencystretch}{3em}
\newcommand{\N}{\mathbb{N}}
\newcommand{\Z}{\mathbb{Z}}
\newcommand{\Q}{\mathbb{Q}}
\newcommand{\R}{\mathbb{R}}
\newcommand{\C}{\mathbb{C}}
\newcommand{\F}{\mathbb{F}}
\newcommand{\A}{\mathbb{A}}
\newcommand{\PP}{\mathbb{P}}
\newcommand{\E}{\mathbb{E}}
\newcommand{\eps}{\varepsilon}
\newcommand{\dd}{\,\mathrm{d}}
\begin{document}
\section*{UnsolvedMath ID: 800585}
\textbf{Source problem code:} catalogue entry
\subsection{Definitions and mathematical statement}
For a set \(X\subseteq\R\), use the subspace topology inherited from the usual topology on \(\R\). Under the standard cover conventions used in the following formulation, a countable Borel cover of \(X\) is an at most countable family \(\mathcal U\) of Borel subsets of \(X\) such that \(\bigcup\mathcal U=X\) and \(X\notin\mathcal U\). Write \(\mathrm B(X)\) for the family of these covers. An \(\omega\)-cover is a cover \(\mathcal U\) such that
\[
  \forall F\subseteq X\text{ finite}\;\exists U\in\mathcal U\quad F\subseteq U.
\]
Write \(\mathrm B_\Omega(X)\) for the countable Borel \(\omega\)-covers of \(X\). A \(\gamma\)-cover is an infinite cover \(\mathcal U\) such that each point of \(X\) belongs to all but finitely many members of \(\mathcal U\); write \(\mathrm B_\Gamma(X)\) for the countable Borel \(\gamma\)-covers.

For families \(\mathcal A\) and \(\mathcal C\) of covers of \(X\), the selection principle \(\mathsf S_1(\mathcal A,\mathcal C)\) means that for every sequence \((\mathcal U_n)_{n\geq1}\) with \(\mathcal U_n\in\mathcal A\), there exist choices \(U_n\in\mathcal U_n\) such that \(\{U_n:n\geq1\}\in\mathcal C\). The principle \(\mathsf S_{\mathrm{fin}}(\mathcal A,\mathcal C)\) means that for every such sequence there are finite, possibly empty, subfamilies \(\mathcal F_n\subseteq\mathcal U_n\) such that
\[
  \bigcup_{n\geq1}\mathcal F_n\in\mathcal C.
\]
Here the union is a family of selected cover members. For a property \(P\) of sets of reals, heredity means
\[
  \forall X\subseteq\R\;\forall Y\subseteq X\quad
  \bigl(P(X)\Longrightarrow P(Y)\bigr).
\]
The subsets \(Y\) are arbitrary; they need not be Borel in \(X\).

The supplied problem asks the following two questions separately:
\begin{enumerate}
\item Is it true that, for every \(X\subseteq\R\) and every \(Y\subseteq X\),
\[
 X\models\mathsf S_1\bigl(\mathrm B_\Omega(X),\mathrm B_\Omega(X)\bigr)
 \quad\Longrightarrow\quad
 Y\models\mathsf S_1\bigl(\mathrm B_\Omega(Y),\mathrm B_\Omega(Y)\bigr)?
\]
\item Is it true that, for every \(X\subseteq\R\) and every \(Y\subseteq X\),
\[
 X\models\mathsf S_{\mathrm{fin}}\bigl(\mathrm B_\Omega(X),\mathrm B_\Omega(X)\bigr)
 \quad\Longrightarrow\quad
 Y\models\mathsf S_{\mathrm{fin}}\bigl(\mathrm B_\Omega(Y),\mathrm B_\Omega(Y)\bigr)?
\]
\end{enumerate}
With the stated convention excluding the whole space from each cover, a finite space has no \(\omega\)-cover, so the two selection properties hold vacuously on finite spaces. The supplied excerpt does not explicitly specify its topology, cover conventions, or selection-principle definitions; this formulation makes the conventional interpretation explicit and requires comparison with the full source.

As context, the supplied excerpt states that none of the selective hypotheses involving open covers is provably hereditary, citing \texttt{[ideals]}. It also states that \(\mathsf S_1(\mathrm B,\mathrm B)\) and all properties of the form \(\Pi(\mathrm B_\Gamma,\mathrm B)\) are hereditary, citing \texttt{[ideals]}, whereas \(\mathsf S_1(\mathrm B_\Omega,\mathrm B_\Gamma)\) is not hereditary, citing \texttt{[MilNonGamma]}. These are supplied contextual claims, not independently verified assertions. The range of \(\Pi\), the scope of the open-cover claim, and the underlying meaning of provability are not specified in the excerpt.
\subsection{Short English statement}
For each of the two stated selection principles, does every subset of a set of reals satisfying that principle also satisfy it? The covers are countable Borel \(\omega\)-covers, and the subsets are arbitrary. The supplied registry labels this problem open; that label and the contextual literature claims have not been independently verified.
\subsection{Sources}
\begin{itemize}
\item Local UnsolvedMath registry, supplied record with permanent ID \(800585\), source problem code \texttt{catalogue entry}, in the AMR Open Problem Lists collection. The record attributes the problem to Tsaban, gives proposed year 2006, and supplies status \texttt{open}. Its upstream revision is \texttt{d371ee917ac9647fbb06cd6fc529deec6f3dda82}.
\item Supplied source attribution: Tsaban, \emph{Selection Principles and Special Sets of Reals: Open Problems} (2006), Problem 11, \url{https://arxiv.org/abs/math/0609601}. The record reports access on 2026-10-16 and extraction method \texttt{arxiv\_latex}; these are supplied provenance details, not a claim of access or verification for this draft.
\item The supplied title is \emph{Selection Principles in Mathematics --- Problem 11 ([ideals, MilNonGamma])}. The excerpt cites \texttt{[ideals]} for its contextual heredity assertions and \texttt{[MilNonGamma]} for its nonheredity assertion. No bibliographic entries identifying those two reference labels are supplied.
\end{itemize}
\end{document}
